\chapter{The Coarsening}
% OUTLINE Ch1 -- the fourth face ENTERS as a COARSENING of the fine argument
% (operator directive 2026-07-26 14:05Z: "volume 6 in chapter 1, described as a
% coarsening of this argument"). THE LOAD-BEARING SCOPING (both chairs, decides
% the chapter's honesty; baked from sentence one): "coarsening" is NOT coarse-
% graining/averaging -- it is an EXACT RE-COUNT (no APPROXIMATION error,
% deterministic, exact-rational) that DEFINES FEWER DISTINCTIONS at the coarser
% floor (defines-not-perceives; not-distinguished != lost; a re-definition not an
% erasure). NEVER "no information lost / all distinctions preserved" unscoped.
% Multigrid/coarse-mesh (Brandt 1977 candidate) as a fenced SHAPE only. Unabashed
% math (presentation law) + four riders; bold-math/fenced-referent; no Lean
% identifiers; step-not-move; lab-gap where any number appears; model-not-world.
% Gate: Kodo (earn/citation) + Beastmaster (arrival-read, "no info lost / all
% preserved" held hardest here).

The fourth face of the model does not begin with a new object. It begins with an
old one --- the fine argument the earlier volumes built, the elaboration and its
count --- read again at a coarser index. That re-reading is the whole mechanism
of this chapter, and before a word of what it produces, the chapter must say
exactly what the re-reading is, because a single loose sentence here would break
the thesis the whole volume rests on.

To coarsen, in this book, is to \emph{re-count}, and re-counting is exact. When
the instrument reads its fine argument at a coarser index it does not average,
does not blur, does not round; there is no approximation in the operation and
nothing is smeared. The coarse count is fixed by the fine one deterministically,
in the exact arithmetic the machine has used since its first volume --- no floats,
no tolerance, no error term. That is the honest and load-bearing difference from
the coarse-graining a working fluid solver performs: a real mesh, restricted to a
coarser grid, \emph{averages} the fine values into the coarse cells and loses the
sub-cell detail as approximation; this instrument does no such thing. What it
computes at the coarse index it computes exactly.

And yet exactness is not the same as keeping every distinction, and the two must
be held apart or the chapter contradicts itself. At the coarser index the
instrument \emph{defines fewer distinctions}. Write $D_{N}$ for the distinctions
the count draws at resolution $N$ --- the pairs it can tell apart with the numbers
it holds there. Coarsening carries the argument from a fine resolution to a
coarser one, and
\[
  D_{\text{coarse}} \;\subseteq\; D_{\text{fine}},
\]
a genuine containment: two things the fine count told apart may fall, at the
coarser floor, into a single box, and there they are one --- decidably, by the
machine's own test, not by any failure of it. What the coarser resolution does
not carry is therefore not \emph{lost}, because nothing was approximated away; it
is \emph{not distinguished}. The machine has defined the two as the same at this
floor, and that is a re-definition, not an erasure. ``Not distinguished'' is not
``gone''; it is ``the same, at this floor.'' Exact in its arithmetic, coarser in
its distinctions: those are the two true things about the re-count, and they do
not conflict, because carrying no approximation error is not the same as
preserving every difference.

\begin{intuition}
The picture worth holding is the multigrid method of scientific computing, where
a problem posed on a fine grid is solved with the help of coarser ones, the coarse
grid capturing the broad shape while the fine grid resolves the detail --- a
technique Achi Brandt set out in 1977. The coarsening here \emph{rhymes} with
that: a fine argument re-read on a coarser index. But the rhyme is a picture and
not the mechanism, and it is fenced where it would mislead. A multigrid
restriction operator averages, and averaging is exactly what this instrument's
re-count does not do. Borrow the image of the coarse grid reading the fine one;
do not borrow the loss. The device's coarsening is the exact half of the picture
without the lossy half.
\end{intuition}

Now the face itself. A fluid is the physicist's name for the behavior that
appears when one stops tracking every particle and watches the collective
instead --- the flow, the dissipation, the shape a great many local interactions
make together.
The model has, from its founding, promised a fluid description of the electron as
one of its four faces, and this is where that description enters: not as an
equation imported from hydrodynamics, but as what the fine argument becomes when
it is re-counted at a coarser floor. The individual steps of the elaboration ---
each exact, each distinguished at the fine resolution --- fall, at the coarser
one, into the boxes the coarse count defines, and what one reads there is no
longer the individual step but the collective structure the steps make. That
collective structure, surviving the re-count, is the model's fluid behavior. The
dissipation, the relaxation, the emergent shape of the later chapters are all
read at this coarser floor, off the same exact argument, seen at the resolution
where the individual falls into the collective.

This is why the volume can call its fourth face a coarsening and mean no
approximation by it. The behavior that emerges is not an averaged summary of the
fine argument, with detail smoothed into a trend. It is exactly what the coarser
floor still tells apart --- the distinctions that survive $D_{\text{coarse}}$,
read off an argument that was never approximated at all. Emergence, in the sense
this book will use the word through every chapter that follows, is precisely
that: not detail smeared into a trend, and not everything carried whole, but the
structure that remains distinguishable when the exact count is re-taken at a
coarser floor.

Two fences close the chapter, as they will close every chapter of this volume.
The fluid behavior described here is the \emph{model's}, an image of flow and not
a fluid of the world; nothing in this coarsening touches the Navier--Stokes
problem of real hydrodynamics, whose name the volume borrows and whose
mathematics it does not. And where the re-count reads a number, that number is the
model's own, at the model's resolution --- never the laboratory's, not here and
not anywhere in these pages. With the coarsening understood exactly, the next
chapter can read the first of the collective structures it reveals: the strain
that relaxes, and the taper that damps it to the floor.
